\begin{partdivider}{I}{LE MONDE AVANT}\end{partdivider}

\bookchap{1}{Nous avons toujours externalisé notre intelligence}

Il est tentant de voir dans l'intelligence artificielle une rupture radicale et sans précédent dans l'histoire de l'humanité. À bien y regarder, cette impression est partiellement justifiée. Elle oublie toutefois un fait beaucoup plus ancien : \textbf{l'humanité n'a jamais pensé seule.}

Depuis ses origines, l'être humain a trouvé le moyen de ne pas tout garder dans sa mémoire, dans son esprit ou dans ses seules capacités physiques. Il a externalisé. Il a délégué à des objets, à des signes ou à des systèmes, certaines tâches qui lui étaient pénibles, coûteuses ou tout simplement impossibles à réaliser seul.

Le premier et sans doute le plus fondamental de ces actes fut le

\textbf{langage}. En donnant un nom aux choses, en les partageant par la parole, nous avons cessé de devoir réinventer l'expérience à chaque individu. Le savoir a pu commencer à circuler au-delà de celui qui l'avait vécu. Déjà là, l'intelligence cessait d'être strictement individuelle pour devenir partiellement sociale et distribuée.

Avec l'apparition de \textbf{l'écriture}, ce mouvement prit une ampleur décisive. Ce qui était dit pouvait survivre à celui qui l'avait prononcé. Ce qui était pensé ne s'évaporait plus avec la voix. Les premières tablettes mésopotamiennes, les hiéroglyphes ou les inscriptions chinoises nous ont permis, pour la première fois, de déléguer notre mémoire à un support matériel. Un homme pouvait, des siècles après sa mort, continuer à transmettre un raisonnement à d'autres hommes qu'il n'aurait jamais rencontrés.

L'invention de \textbf{l'imprimerie} poussa cette logique beaucoup plus loin. Gutenberg ne révolutionna pas seulement la diffusion des livres. Il rendit possible la multiplication à grande échelle d'idées stabilisées. Là où un manuscrit était rare et précieux, un ouvrage imprimé pouvait se répandre. L'accès à une partie du savoir cessa peu à peu d'être l'apanage d'une élite restreinte de clercs ou de lettrés. L'intelligence écrite sortit de ses cercles fermés.

À mesure que les sociétés se complexifiaient, cette externalisation s'étendit à d'autres domaines. Les \textbf{cartes géographiques} déléguèrent à un dessin notre capacité à nous représenter un espace immense que nous ne pouvions pas embrasser du regard. Les \textbf{tables de logarithmes}, puis les

\textbf{calculatrices mécaniques}, allégèrent notre effort de calcul. Pendant des siècles, il exista même des «calculateurs humains» spécialisés, dont le travail finit par être relayé par des machines dès que celles-ci devinrent assez fiables.

Avec \textbf{Internet}, ce mouvement atteignit un nouveau palier. Pour la première fois dans l'histoire de l'humanité, l'essentiel des savoirs accumulés devint potentiellement interrogeable à distance. Les moteurs de recherche nous ont habitués à ne plus mémoriser certains faits : adresses, dates, formules, définitions ou itinéraires. Nombre d'entre nous avons cessé de retenir des numéros de téléphone, parce que notre téléphone s'en charge à notre place. Nous nous fions au GPS plutôt qu'à notre orientation dans l'espace. Chacun de ces petits renoncements, presque imperceptibles au quotidien, illustre la même logique : nous déléguons pour libérer notre attention ailleurs.

Ce phénomène n'est pas nouveau. Il a même été théorisé bien avant l'IA. Le psychologue soviétique Lev Vygotsky insistait sur le fait que notre pensée ne se développe pas à l'intérieur d'un individu isolé, mais par l'intermédiaire d'outils et dans l'interaction avec autrui. Plus près de nous, Andy Clark et David Chalmers ont proposé l'idée de l'«esprit étendu» (\emph{Extended Mind}). Selon eux, certains objets extérieurs — un carnet, un agenda, un ordinateur — ne se contentent pas d'aider notre cognition : ils en font partie intégrante lorsque nous les utilisons de manière systématique pour résoudre un problème.

Il faut donc l'admettre avec réalisme : l'être humain est, par nature, un être cognitif distribué. Il pense avec le langage, avec l'écriture, avec ses outils, avec les autres. Il n'a jamais été un esprit entièrement clos sur lui-même.

Cela nous conduit à la question centrale de ce chapitre :

\begin{aiquote}
L'intelligence artificielle est-elle simplement la dernière étape logique de cette très longue histoire d'externalisation de nos capacités cognitives, ou constitue-t-elle un changement de nature ? À première vue, la continuité est frappante. Comme l'écriture nous a déchargés de la mémoire, comme Internet nous a déchargés de la recherche d'informations, l'IA semble nous décharger d'une part croissante de la formulation, de l'analyse ou de la synthèse. Sur ce plan, il existe une incontestable continuité historique.
\end{aiquote}

Pourtant, trois différences méritent d'être soulignées avec prudence, en les présentant davantage comme des \textbf{observations} que comme des affirmations définitives.

La première différence tient à \textbf{ce qui est externalisé.} Les technologies précédentes externalisaient principalement des fonctions bien délimitées : stocker, calculer, localiser ou retrouver. L'IA générative, elle, est capable d'externaliser non plus un résultat figé, mais un processus : formuler un raisonnement, proposer des hypothèses, structurer une argumentation ou adapter un discours à un contexte donné.

La deuxième différence est d'ordre \textbf{d'interactivité.} Un livre ne répond pas lorsqu'on lui pose une question. Un moteur de recherche nous renvoie des sources, il ne dialogue pas avec notre intention. Une IA conversationnelle, à l'inverse, s'adapte au fil de l'échange. Elle peut reformuler, préciser, nuancer à notre demande. Ce caractère dialogique crée un rapport nouveau entre l'utilisateur et l'outil.

La troisième différence concerne \textbf{l'accessibilité et la généralisation.} Jamais auparavant une capacité cognitive aussi polyvalente n'a été rendue disponible, à un coût très faible, à des milliards de personnes, sans nécessiter de formation technique particulière pour l'utiliser.

L'hypothèse qui se dessine est la suivante : ces différences sont suffisantes pour parler d'un changement d'échelle, et peut-être, dans certains domaines, d'un changement de nature. Il serait toutefois excessif, à ce stade de son développement, de conclure définitivement à un changement de nature dans tous les domaines de la vie sociale.

Le point important à retenir pour la suite de ce livre est le suivant : nous ne partons pas de zéro. L'humanité a une longue expérience de l'externalisation de son intelligence. À chaque fois, cette délégation nous a permis de libérer des énergies nouvelles, mais elle nous a aussi coûté quelque chose : une partie de nos savoirs, de nos automatismes ou de notre mémoire s'est affaiblie au profit de l'outil.

La question qui se profile déjà est donc double : Qu'allons-nous gagner cette fois-ci en déléguant davantage ? Et qu'est-il légitime, ou souhaitable, de ne pas déléguer ? Nous comprendrons mieux cette inquiétude en interrogeant un autre aspect fondamental de nos sociétés : pourquoi, pendant des siècles, avons-nous attribué tant de valeur à l'intelligence elle-même ?

\bookchap{2}{La valeur de l'intelligence}

Pour saisir ce qui pourrait se trouver bouleversé avec l'IA, il ne suffit pas de regarder son évolution technique. Il faut porter un regard historique sur nos sociétés. Car ce n'est pas un hasard si, dans la plupart des civilisations modernes, l'intelligence, les diplômes, l'expertise ou les qualifications ont été si fortement valorisés sur le marché du travail et dans la hiérarchie sociale.

Pendant très longtemps, certaines formes d'intelligence ont été rares. Un bon calculateur, un juriste compétent, un médecin formé, un ingénieur ou un traducteur ne se formaient qu'au prix d'années d'études. Leur temps était précieux, leur savoir difficilement reproductible. Cette rareté a naturellement trouvé son expression économique : on payait pour accéder à leur jugement, à leur analyse ou à leur capacité à résoudre un problème complexe.

Il faut l'observer avec clarté : \textbf{dans une économie de la rareté, ce qui est difficile à produire, à acquérir ou à reproduire tend à prendre de la valeur.} L'expertise intellectuelle a fonctionné selon ce principe pendant des siècles.

Ce lien entre intelligence, compétence, travail, revenu et statut social s'est progressivement renforcé avec la modernité. L'industrialisation valorisait d'abord la force de travail et la maîtrise des machines. Au fil du XX\up{e} siècle, les économies se sont de plus en plus tournées vers les services et vers ce que l'on a appelé l'«économie de la connaissance». On a alors commencé à mesurer la progression sociale par l'accès à des formations de plus en plus longues, dans l'espoir légitime que l'accumulation de savoir garantirait une meilleure insertion professionnelle et une forme de mobilité sociale.

Ce système reposait sur un postulat implicite, rarement explicité, mais extrêmement solide : \textbf{l'intelligence appliquée à un problème constitue une ressource rare, et cette rareté justifie sa valorisation économique et sociale.}

Aujourd'hui, ce postulat commence à être interrogé. Il serait cependant prématuré de l'affirmer comme caduc. Il convient plutôt de formuler une

\textbf{hypothèse} : si certaines compétences intellectuelles jusqu'ici rares

deviennent beaucoup plus faciles à reproduire, à mettre à disposition et à accéder à grande échelle, alors le fondement même de leur valeur peut se déplacer.

Pour mieux comprendre ce déplacement potentiel, prenons un exemple simple et réaliste.

Autrefois, dans certaines entreprises, il existait des «calculateurs humains». Leur rôle était d'effectuer des opérations longues et fastidieuses. Lorsqu'apparurent les premières calculatrices fiables, puis les ordinateurs, ce métier ne disparut pas brutalement dans tous les secteurs, mais la valeur économique attachée à la simple capacité de calcul chuta très fortement. Ce qui était jadis un savoir valorisé devint progressivement une commodité attendue. L'humain n'y perdait pas nécessairement toute utilité, mais ce qui justifiait qu'on le paie spécifiquement pour cette tâche s'était érodé.

Ce cas est instructif. Il nous permet d'établir une distinction importante, que nous retrouverons tout au long de ce livre :

C'est un fait établi : une tâche peut être automatisée.

On observe ensuite que lorsque cette tâche cesse d'être rare, la valeur qui lui était attachée a tendance à se déplacer vers d'autres éléments (le contexte, l'interprétation, la prise de décision ou la responsabilité).

Il serait abusif de généraliser mécaniquement cet exemple à l'ensemble des activités intellectuelles. Un raisonnement juridique, une analyse stratégique, une interprétation historique ou un jugement éthique sont bien plus complexes qu'un simple calcul. Ils impliquent du contexte, des nuances, des incertitudes et des finalités qui ne se réduisent pas à l'exécution.

Néanmoins, l'IA générative montre qu'une partie non négligeable de ce qui était autrefois considéré comme du «travail intellectuel de routine» peut aujourd'hui être produit, reformulé, synthétisé ou structuré avec une aide artificielle. Cela soulève une question inévitable :

\begin{aiquote}
Que se passe-t-il lorsque des compétences qui ont fondé, pendant des générations, notre idée même de la réussite sociale — savoir analyser, rédiger, raisonner, coder — cessent, pour une part significative, d'être des marqueurs distinctifs ?
\end{aiquote}

L'interprétation qui s'impose est la suivante : historiquement, nos sociétés ont récompensé l'effort nécessaire pour acquérir une compétence rare. Si cet effort d'accès se trouve considérablement réduit, il est logique de penser que la valeur se déplace. Le problème n'est pas de savoir si elle disparaît, mais de se demander \emph{vers quoi} elle se déplace.

Trois pistes méritent d'être explorées, sans pour autant les présenter comme certaines :

Premièrement, la valeur pourrait se déplacer vers \textbf{le jugement}. Être capable de formuler une bonne réponse est une chose ; être capable de juger si cette réponse est pertinente, juste, utile ou éthique dans une situation donnée en est une autre.

Deuxièmement, elle pourrait se déplacer vers \textbf{la connaissance du problème, et non plus vers la connaissance de la solution.} Celui qui sait le mieux poser les bonnes questions, cerner un enjeu réel ou définir ce qui est réellement important dans un contexte précis conserve un avantage humain difficile à automatiser.

Troisièmement, elle pourrait se déplacer vers \textbf{les dimensions relationnelles.} Dans de nombreux métiers (santé, éducation, accompagnement, conseil), ce n'est pas uniquement l'expertise qui est recherchée, c'est la confiance, l'écoute, la présence ou la capacité à prendre en compte l'autre dans toute sa complexité.

Reste une hypothèse prudente : Lorsque l'intelligence devient moins rare, ce qui reste difficilement reproductible par une machine — le contexte singulier, le jugement situé, la responsabilité ou la relation humaine — a des chances de prendre un poids nouveau dans la manière dont nous attribuons de la valeur au travail.

Il faut cependant se garder d'un optimisme trop hâtif. Rien ne garantit que cette valeur nouvelle sera reconnue, rémunérée ou équitablement répartie dans nos économies actuelles. C'est là tout le paradoxe.

Pendant des siècles, nous avons fondé une grande partie de notre organisation sociale sur la rareté de l'intelligence. Si cette rareté venait à s'atténuer à une échelle sans précédent, nous serions alors confrontés à une possibilité historique tout à fait inédite.

C'est ce que nous allons explorer dans le chapitre suivant : et si l'intelligence devenait, pour la première fois, une ressource abondante ?

\bookchap{3}{Une nouvelle abondance}

Jusqu'à présent, nous avons constaté deux choses. D'une part, l'humanité a toujours eu recours à des outils pour externaliser une partie de son intelligence. D'autre part, nos sociétés ont structuré la valeur économique autour de la rareté de certaines compétences intellectuelles. Il existe donc un lien étroit entre rareté de l'intelligence et organisation de notre société.

Le chapitre précédent nous a laissé une question en suspens : que se produit-il lorsque cette rareté commence à diminuer à grande échelle ?

Nous pouvons formuler une hypothèse simple :

\begin{aiquote}
Et si, pour la première fois dans l'histoire, l'intelligence cessait d'être un facteur limitant pour devenir une ressource irrigant toute l'économie ?
\end{aiquote}

Il faut prendre soin de bien définir ce que nous entendons ici par «abondance». Nous ne prétendons pas que toute forme d'intelligence deviendra gratuite ou banale. Il s'agit d'une \textbf{hypothèse d'évolution}, limitée à certaines capacités cognitives : expliquer, synthétiser, analyser, rédiger, traduire, programmer, raisonner sur un domaine connu ou structurer un problème.

Pendant des générations, de nombreuses sociétés se sont heurtées à des pénuries récurrentes : manque d'enseignants qualifiés dans certains territoires, manque d'ingénieurs, de médecins, d'experts agricoles, d'analystes ou encore de développeurs de logiciels. Ces pénuries limitaient leur capacité à se développer, à moderniser leurs services ou à réduire les inégalités d'accès au savoir.

On observe que l'une des conséquences potentielles de l'IA est de rendre certaines de ces expertises beaucoup plus facilement accessibles, y compris là où les ressources humaines qualifiées faisaient historiquement défaut. Cette possibilité mérite d'être soulignée, notamment dans une perspective mondiale.

À titre d'exemple réaliste : un élève vivant dans une zone reculée, qui n'a pas facilement accès à un professeur particulier de mathématiques ou de physique, peut aujourd'hui, avec les précautions nécessaires, s'appuyer sur un tuteur IA pour s'expliquer un concept à son rythme. De même, un petit exploitant agricole dans certaines régions d'Afrique peut, à condition d'avoir un accès minimal à Internet, obtenir des explications vulgarisées sur des pratiques culturales adaptées à son contexte local. Il s'agit ici d'un possible, et non d'une généralisation. Cette idée d'une forme d'«intelligence abondante» est séduisante. Elle fait naître l'espoir légitime d'une démocratisation partielle de l'expertise. Pourtant, il faut immédiatement lui opposer un important principe de prudence : \textbf{l'abondance d'une ressource ne produit pas automatiquement une société plus juste, plus sage ou plus prospère.}

L'histoire nous fournit plusieurs mises en garde à cet égard. Prenons l'exemple de l'information. Avec Internet, nous sommes passés d'une relative rareté de l'information à une véritable explosion de celle-ci. Très peu de personnes prédisaient, à ses débuts, qu'au lieu de nous rapprocher d'une société mieux informée, nous nous retrouverions aussi confrontés à la désinformation, à la surcharge informationnelle, à la polarisation ou à la difficulté grandissante de distinguer ce qui est fiable de ce qui ne l'est pas.

Autre exemple : l'abondance alimentaire dans de nombreux pays développés n'a pas supprimé la faim dans le monde. Elle a coexisté avec des problèmes nouveaux : obésité, gaspillage et inégalités d'accès criantes à cette même abondance. Le problème n'était donc pas la production, mais sa répartition et son organisation sociale.

Ces analogies nous invitent à distinguer clairement deux niveaux :

\begin{eplist}
\item \textbf{Le niveau technique :} Peut-on rendre une capacité cognitive plus abondante ? Sur certains pans, la réponse est : probablement oui.
\item \textbf{Le niveau social, économique et politique :} Qui bénéficiera réellement de cette abondance ? Comment sera-t-elle distribuée ? À quelles fins sera-t-elle utilisée ?
\end{eplist}
L'hypothèse qui se dessine est la suivante : le véritable enjeu de l'IA ne se situe pas dans sa capacité à produire de l'intelligence à grande échelle, mais dans la manière dont les bénéfices de cette abondance potentielle seront partagés au sein des sociétés.

Un autre paradoxe mérite d'être évoqué. En économie, on connaît le «paradoxe de Jevons». Au XIXe siècle, William Stanley Jevons observa que lorsque l'efficacité dans l'utilisation d'une ressource (le charbon) s'améliorait et que son coût baissait, la consommation totale de cette ressource ne diminuait pas : elle augmentait. Autrement dit, plus un bien devient moins cher à utiliser, plus on a tendance à l'utiliser massivement, ce qui peut compenser, voire dépasser les gains d'efficacité.

Nous pouvons, par simple analogie et à titre d'\textbf{interprétation}, nous demander s'il n'en irait pas de même avec l'intelligence. Si l'analyse, la rédaction, la programmation ou la création deviennent beaucoup moins coûteuses à produire grâce à l'IA, il est très plausible que la demande pour ces activités explose au lieu de se réduire. On demandera davantage d'analyses, davantage de contenus, davantage de logiciels, davantage de personnalisations. L'intelligence consommée pourrait croître considérablement, tout en modifiant profondément ce que nous attendons d'elle.

Ce paradoxe nous met en garde contre une erreur de raisonnement trop linéaire : l'abondance d'une capacité ne signifie pas nécessairement la baisse de son usage, elle peut au contraire en multiplier les applications.

Il convient également de ne pas confondre deux notions différentes :

\textbf{abondance} et \textbf{sagesse}. Une société peut disposer d'une quantité jamais atteinte d'explications, d'analyses ou d'informations et, en même temps, éprouver d'autant plus de difficultés à discerner ce qui est pertinent, ce qui est juste ou ce qui importe véritablement. L'intelligence disponible à grande échelle n'équivaut en aucun cas à une meilleure capacité collective de jugement.

Ce que nous retiendrons de ce chapitre : À titre d'hypothèse de travail pour la suite du livre, nous retiendrons l'idée suivante : l'IA pourrait faire basculer certaines formes d'intelligence d'une économie fondée sur la rareté vers une économie où elles tendent à devenir plus abondantes. Ce basculement historique est potentiellement profond. Toutefois, ni l'abondance technique ni l'accessibilité ne suffisent à garantir une société plus équitable. La question sociale demeure entière : \emph{Que ferons-nous de cette nouvelle abondance si elle se concrétise ?}

Le mouvement historique que nous venons de décrire — cette longue habitude qu'a l'humanité d'externaliser son intelligence — nous aide à mieux mesurer l'ampleur du phénomène. Il nous reste maintenant à comprendre comment, techniquement et conceptuellement, nous sommes passés d'un simple outil numérique à un système avec lequel il devient possible de penser, dans une certaine mesure, \emph{avec} nous. C'est l'objet de la Partie II : Le Grand Basculement.
